
\begin{proposition}[Complete primary and secondary nonlinear Hopf star criterion]\label{cdp:st:C.2.UN.50}
In \ref{cdp:st:C.2.UN.49} allow each prescribed leaf to be
\[
 H_i=H_{r_i}(\beta_i)/\mu_{n_i},\qquad
 \zeta\cdot(z,w)=(\zeta^{r_i}z,\zeta w),\qquad
                         \gcd(r_i,n_i)=1,               \tag{C.SN.1}\label{cdp:eq:C.SN.1}
\]
including $n_i=1$. These are all nonlinear primary and
secondary Hopf surfaces of algebraic dimension zero,
by \ref{cdp:st:C.2.UN.43}. The same equivalence between an actual
smooth-total torus smoothing and compact realization
over $\PP^1$ holds. Its necessary-and-sufficient tests
are dual normal lines, an elliptic central base with
no unmarked multiple fibres, one common resonance
$r_i=r$, and zero integral Chern vector of the
inverse logarithmic transform defined below.

The marked full fibres in $P$ still have multiplicity
$r_i$, while the finite integer $n_i$ enters the
exact inverse-transform framing. In a global
punctured-base trivialization its meridian is
\[
       \mu_{P,i}=r_i\delta_{\rm ext}
                         +(r_i\nu_i-n_i e_i),\qquad
                         c(T_0)=\sum_i\nu_i.             \tag{C.SN.2}\label{cdp:eq:C.SN.2}
\]
Thus its rational Seifert vector is
$\sum_i(\nu_i-n_i e_i/r_i)$ in this convention.
The full integer $n_i$ is retained; its residue
modulo $r_i$ alone does not determine the criterion.
All original component moduli and elliptic gluing
maps are preserved by the compact realization.
\end{proposition}

\begin{proof}
We give every change to the exact maps of \ref{cdp:st:C.2.UN.49}.
Fix one leaf and suppress its index. Choose integers
$k,l$ with
\[
                         rk-nl=1,\qquad \gamma=\beta^n.
                                                        \tag{C.SN.3}\label{cdp:eq:C.SN.3}
\]
On the primary cover of the paired $P$ neighborhood
use $p=x^rz$, the infinite deck
$(z,x)\mapsto(\beta^rz,\beta^{-1}x)$, and the finite
action $(z,x)\mapsto(\zeta^rz,\zeta^{-1}x)$.
The following are genuine invariant coarse coordinates:
\[
 u=z^n,\qquad X=xz^k,\qquad
 p=X^ru^{-l},\qquad
 (u,X)\sim(\gamma^ru,\gamma^lX).
                                                        \tag{C.SN.4}\label{cdp:eq:C.SN.4}
\]
Here $u\ne0$. Since $rk-1=nl$, both coordinates
are finite-group invariant, and the finite quotient
map has no ramification. Conversely, a root of $u$
recovers $z$ and $x=Xz^{-k}$; all choices are exactly
the original finite action. Thus \eqref{cdp:eq:C.SN.4} is the
actual coarse neighborhood, not a formal coordinate
substitution.

The normal of $J$ in $P$ has multiplier $\gamma^l$
over $\C^*/\gamma^{r\Z}$. Its order is exactly $r$,
since $\gcd(l,r)=1$. On the normal root cover write
$u=y^r$, $X=s y^l$; then $p=s^r$.
The generator $(y,s)\mapsto(\zeta y,\zeta^{-l}s)$
preserves these coordinates. Normalize its normal
character to $\zeta$: the fibre translation is then
$y\mapsto\zeta^n y$, because $-ln\equiv1\pmod r$.
This proves the exact torsion-character comparison,
including its dependence on $n$.
The dual normal on the secondary Hopf axis is
the inverse character, so its normal-root
identification with the given $P$ gives precisely
\eqref{cdp:eq:C.SN.4}.

The punctured product map used for the inverse
transform is
\[
 (u,X)\longmapsto(p,W)
       =(X^ru^{-l},\,X^{-n}u^k),\qquad
 u=p^nW^r,\quad X=p^kW^l,\qquad
                       F=\C^*/\gamma^{\Z}.               \tag{C.SN.5}\label{cdp:eq:C.SN.5}
\]
The exponent matrix has determinant $rk-nl=1$.
Consequently these are inverse holomorphic maps
on the punctured chart. On the primary cover
they are $W=x^{-n}$ and $u=z^n=p^nW^r$.
Replacing the chart by $F\times\Delta_p$ through
\eqref{cdp:eq:C.SN.5} defines the actual principal torsor $T_0$.
It retains the original analytic transitions
exactly as in (C.NS.3).

Let $a$ be the primitive positive angular loop of
$u$ with $X$ fixed, $b$ the infinite deck loop,
and $\mu_P$ the normal loop of $X$. In the
punctured product basis $(\delta,e,h)$,
\eqref{cdp:eq:C.SN.5} gives the full integral matrix
\[
 \begin{pmatrix}a&b&\mu_P\end{pmatrix}_{(\delta,e,h)}
   =
 \begin{pmatrix}
   -l&0&r\\
    k&0&-n\\
    0&1&0
 \end{pmatrix},
 \qquad
 \delta=n a+k\mu_P,\quad e=r a+l\mu_P,\quad h=b.
                                                        \tag{C.SN.6}\label{cdp:eq:C.SN.6}
\]
Its determinant is one. In particular the normal
meridian is $r\delta-ne$ and is primitive even
when $n>1$.

On the Hopf side the invariant normal coordinate
is $Y=wz^{-k}$. Its normal-line multiplier is $\gamma^{-l}$,
and the diagonal reference pair has reciprocal
base coordinate $Y^ru^l=w^r/z$ and fibre coordinate
$w^n=Y^nu^k$. The complementary solid torus
therefore has meridian $n a-k\mu_H$ and fibre
cycles $(r a-l\mu_H,b)$. The actual smoothing
identifies $\mu_H=-\mu_P$. Formula \eqref{cdp:eq:C.SN.6}
then identifies its meridian with $\delta$
and its fibre cycles with $(e,h)$. This proves
the inverse-transform topology and the common
integral fibre-cycle identification for every
secondary leaf. Substituting
$\delta_i=\delta_{\rm ext}+\nu_i$ proves \eqref{cdp:eq:C.SN.2}.
The nonlinear-to-diagonal comparison (C.NS.5)
is finite-group invariant, so its boundary
calculation applies to this actual quotient.

The canonical form
$dz\wedge dw/w^{r+1}$ is invariant under both
the infinite contraction and $\mu_n$. It gives
$K_H=-(r+1)J$ on the quotient. Thus the
canonical-order calculation, the effectivity
argument forcing elliptic $B$ and no unmarked
multiple fibres, and the Chern-zero necessity
are unchanged. The only change to the core
period is explicit. In $U=z/w^r$, $V=\log w$,
the two primitive periods are
\[
             (0,2\pi i/n),\qquad
             (\beta^{-r},\log\beta).
                                                        \tag{C.SN.7}\label{cdp:eq:C.SN.7}
\]
The leading two-form has core integral
$\pm2\pi i c\beta^{-r}/n$, still nonzero.
All cores are the same integral fibre class
by \eqref{cdp:eq:C.SN.6}, so their one global period
again forces all resonances $r_i$ to coincide.
The finite orders $n_i$ need not coincide.

For compact sufficiency perform the exact
inverse transformations and use their
Chern-zero torsor $T_0$, which is an actual
torus by (C.NS.14). In the twistor family
(C.NS.15), choose the $i$-th fibre basis
so that its multiplicative modulus is
$\gamma_i=\beta_i^{n_i}$. After $t=s^r$,
the same four-period calculation gives
\[
 \mathcal T_s=M_{i,s}^{\times}/
   \langle(p,\ell)\mapsto
        (p+d_i s^r,\gamma_i\ell)\rangle,\qquad d_i\ne0.
                                                        \tag{C.SN.8}\label{cdp:eq:C.SN.8}
\]
At the marked point perform the weight-$(r,1)$
blowup, but now tensor the flat line by
$\OO(n_iR_i)$. Its local line coordinate
has weight $-n_i$. Setting it equal to one
leaves the finite action
\[
 (z,w,x)\mapsto
      (\zeta^{r}z,\zeta w,\zeta^{-1}x),
                    \qquad \zeta^{n_i}=1.                \tag{C.SN.9}\label{cdp:eq:C.SN.9}
\]
It is free: a nonzero $w$ or $x$ forces
$\zeta=1$, and if both vanish then $z\ne0$
and $\gcd(r,n_i)=1$ forces the same conclusion.
This proves directly that the coarse nonzero
line is smooth. Equivalently the weighted
orbifold inertia has order $r$ and acts
on the line with the faithful weight $n_i$.

Choose the given $n_i$-th root $\beta_i$
of $\gamma_i$. The infinite line deck
lifts to the exact map
\[
 (z,w,x)\mapsto
   \bigl(\beta_i^r(z+d_iw^r),\beta_iw,\beta_i^{-1}x\bigr)
                                                        \tag{C.SN.10}\label{cdp:eq:C.SN.10}
\]
on \eqref{cdp:eq:C.SN.9}. It commutes with the finite
action. Its exceptional component is
exactly $H_r(\beta_i)/\mu_{n_i}$ after
$Z=z/(\beta_i^r d_i)$.
Off that component the flat-line
coordinate is $W=x^{-n_i}$; consequently
the remaining central chart and its
outer product are exactly
\eqref{cdp:eq:C.SN.4}--\eqref{cdp:eq:C.SN.5}.

The polynomial norm distortion is now
$C(1+|j|)^{n_i/r}$ for the $j$-th
translation lift. The flat-line
contraction multiplier is
$|\gamma_i|\exp(-\operatorname{Re}(c_i(s)d_i s^r))$,
uniformly less than one on a small disc.
It still dominates that polynomial.
The same compact norm annulus proves
properness, and comparison with arbitrarily
large powers proves freeness of the
infinite cyclic quotient. Thus there
is again an actual isomorphism of the
complement of the new Hopf component
with the original twistor tube minus
its marked central elliptic curve.

In the cluster calculation the original
flat coordinate is $x^{-n_i}$, so only
$O(1+|\log|x||)$ deck iterates are
needed. Their base displacement remains
$O(|s|^r(1+|\log|x||))$. All cluster
points therefore lie on the prescribed
central fibre curve. The simultaneous
disjoint replacements and compact gluing
of \ref{cdp:st:C.2.UN.49} apply with these exact
coarse maps. The main component is the
original $P$, and the finite diagonal
actions commute with the axis-translation
automorphisms used there. Every original
gluing map is retained. The full
central equations are still $xw=s$;
the finite action preserves this
function and is free, so the full
fibre remains reduced SNC.
This proves the equivalence and
compact realization asserted here.
\end{proof}


